\section*{Hoofdstuk \ref{ch:b1:rifts-ridges-margins} --- Slenken, ruggen en passieve randen}

\begin{solution}{exo:b1:rifts-ridges-margins:1}
(a) De gang zou $\qty{60000}{m}\times\qty{8}{m} = \qty{4.8e5}{m^2}$ in grondplan beslaan, dus haar gemiddelde hoogte zou $\qty{2.5e9}{m^3}/\qty{4.8e5}{m^2} =
\qty{5200}{m}$ zijn, ongeveer \qty{5}{km}: de gang sneed door een groot deel van de \omterm{def:b1:layered-earth:layers}{korst}.
(b) $\qty{8}{m}/\qty{0.018}{m.yr^{-1}} = 444$ jaar, ongeveer vier eeuwen.
\end{solution}

\begin{solution}{exo:b1:rifts-ridges-margins:2}
(a) De \omterm{def:b1:cross-sections:throw}{heave} is de \omterm{def:b1:cross-sections:throw}{worp} gedeeld door de tangens van de helling:
$3.0/\tan\ang{60} = \qty{1.73}{km}$ en $0.8/\tan\ang{70} = \qty{0.29}{km}$;
de uitrekking van de doorsnede is \qty{2.02}{km}. (b) Vóór de breuken was de
doorsnede $25 - 2.02 = \qty{22.98}{km}$ breed, ongeveer \qty{23}{km}: zij is met
ongeveer \qty{9}{\percent} verlengd.
\end{solution}

\begin{solution}{exo:b1:rifts-ridges-margins:3}
Ultralangzaam (onder \qty{20}{mm.yr^{-1}}): de Gakkelrug (12), de centrale Rode Zee (15) en de Reykjanesrug (19), de laatste dicht bij de grens. Langzaam (20 tot \qty{55}{mm.yr^{-1}}): de westelijke Carlsbergrug (23) en de Zuid-Atlantische Oceaan (33). Snel (boven \qty{75}{mm.yr^{-1}}): de Oost-Pacifische Rug (143). De ultralangzame en langzame ruggen zouden een as-dal tussen normale breuken moeten hebben; de Oost-Pacifische Rug een as-hoogte.
\end{solution}

\begin{solution}{exo:b1:rifts-ridges-margins:4}
(a) $\qty{3.0}{km^2}\times\qty{7.1}{km} = \qty{21.3}{km^3}$, ongeveer
\qty{21}{km^3} per jaar. (b) $\qty{3.028e8}{km^2}/\qty{3.0}{km^2.yr^{-1}} =
\qty{1.01e8}{yr}$, \qty{101}{Myr}. De oudste bodem, ongeveer \qty{190}{Myr}, is ouder omdat de bodem niet in volgorde van leeftijd wordt vernietigd: subductie verbruikt bodem van elke leeftijd, sommige oude bodem overleeft waar geen trog haar heeft bereikt, en het productietempo heeft gevarieerd.
\end{solution}

\begin{solution}{exo:b1:rifts-ridges-margins:5}
(a) $2500 + 350\times2 = \qty{3200}{m}$, $2500 + 350\times6 = \qty{4600}{m}$,
$2500 + 350\times8 = \qty{5300}{m}$. (b) $t = \bigl((d - 2500)/350\bigr)^2$:
$(2100/350)^2 = \qty{36}{Myr}$ en $(2800/350)^2 = \qty{64}{Myr}$. (c) De bodem zakt
\qty{1400}{m} tussen 4 en \qty{36}{Myr}, maar slechts \qty{700}{m} tussen 36 en
\qty{64}{Myr}, over een vergelijkbare tijd.
\end{solution}

\begin{solution}{exo:b1:rifts-ridges-margins:6}
Het effectieve uitzettingscoëfficiënt is
$\num{3.2e-5}\times3300/2300 = \qty{4.591e-5}{K^{-1}}$, en
$\sqrt{\kappa\times\qty{1}{Myr}} = \sqrt{\num{e-6}\times\num{3.156e13}} =
\qty{5618}{m}$. Dan $C = (2/\sqrt\pi)\times\num{4.591e-5}\times1350\times5618 =
1.1284\times\num{4.591e-5}\times1350\times5618 = \qty{393}{m.Myr^{-1/2}}$. Voor een
\omterm{def:b1:layered-earth:layers}{mantel} die \qty{250}{K} kouder is, schaalt $C$ met $T_m - T_0$:
$393\times1100/1350 = \qty{320}{m.Myr^{-1/2}}$. De bodem van een kouderen oceaan
zakt trager en volgt de kleinere constante. De gefitte \num{350} van het
hoofdstuk is zo'n geval niet: zij ligt binnen de onzekerheid van de parameters.
\end{solution}

\begin{solution}{exo:b1:rifts-ridges-margins:7}
(a) $\qty{e-6}{cal.cm^{-2}.s^{-1}} = \num{4.184e-6}\times\num{e4} =
\qty{4.184e-2}{W.m^{-2}} = \qty{41.84}{mW.m^{-2}}$, dus de constante is
$11.3\times41.84 = \qty{473}{mW.m^{-2}.Myr^{1/2}}$, \qty{3.5}{\percent} lager dan de
490 van \cref{ch:b1:heat-budget}: de twee aanpassingen komen overeen. (b) $q = 473$,
$473/\sqrt{10} = 150$ en $473/\sqrt{150} = \qty{38.6}{mW.m^{-2}}$. De wet van de
vierkantswortel blijft dalen, terwijl het plaatmodel tot een constante afvlakt: bij
\qty{150}{Myr} is de waarde van de wet van de vierkantswortel te laag, zoals de
diepte-wet te diep gaat.
\end{solution}

\begin{solution}{exo:b1:rifts-ridges-margins:8}
Neem het noordelijke segment op $x = 0$ en het zuidelijke op $x = -40$
(km), de transform op $y = 0$ ertussen. (a) Net ten noorden van de transform ligt
de bodem ten westen van het noordelijke segment en beweegt naar het westen; net
ten zuiden ervan ligt zij ten oosten van het zuidelijke segment en beweegt naar het
oosten. Gezien vanaf beide zijden gaat de verre zijde naar links: de verschuiving is
linkszijdig (sinistraal). (b) De rug lijkt rechtszijdig verplaatst, met haar
zuidelijke deel naar het westen: de tegengestelde zin, zoals
\cref{prop:b1:rifts-ridges-margins:transform} voorspelt. (c) Aardbevingen komen
alleen voor op de \qty{40}{km} tussen de twee uiteinden van de rug, waar de twee
zijden bij verschillende platen horen; verder, langs de \omterm{def:b1:rifts-ridges-margins:transform}{breukzones}, horen beide
zijden bij dezelfde plaat en bewegen zij samen.
\end{solution}

\begin{solution}{exo:b1:rifts-ridges-margins:9}
De halve snelheid is \qty{10.6}{mm.yr^{-1}}, \qty{10.6}{km.Myr^{-1}}. (a) De leeftijdssprong is $350/10.6 = \qty{33.0}{Myr}$. (b) Aan het einde van de transform staat bodem van \qty{33.0}{Myr} oud ($2500 + 350\sqrt{33.0} = \qty{4511}{m}$) tegenover bodem van leeftijd 0 (\qty{2500}{m}): een sprong van ongeveer \qty{2010}{m}. \qty{100}{km} verder zijn de leeftijden $450/10.6 = 42.5$ en $100/10.6 = \qty{9.4}{Myr}$, de diepten 4781 en \qty{3575}{m}: een sprong van ongeveer \qty{1210}{m}. \qty{1000}{km} verder zijn de leeftijden 127.4 en \qty{94.3}{Myr}, de diepten 6450 en \qty{5899}{m}: \qty{550}{m}. (Deze leeftijden liggen boven \qty{80}{Myr}, waar de wet van de vierkantswortel te diep gaat; de plaatwet zou een kleinere sprong geven.)
\end{solution}

\begin{solution}{exo:b1:rifts-ridges-margins:10}
(a) Met gelijke oppervlakten voor elke leeftijd is de gemiddelde diepte het gemiddelde van $d(t)$ over $[0, A]$:
\[
  \frac1A\int_0^A\bigl(d_r + C\sqrt t\bigr)\mathrm dt = d_r + \frac CA\cdot\frac23A^{3/2}
  = d_r + \frac23C\sqrt A.
\]
(b) Voor $A = \qty{180}{Myr}$: $2500 + \frac23\times350\times13.416 =
\qty{5630}{m}$; voor \qty{120}{Myr}: $2500 + \frac23\times350\times10.954 =
\qty{5056}{m}$, \qty{574}{m} ondieper. (c) Laat de zee met $r$ stijgen en de belaste bodem met $s = r\rho_w/(\rho_m - \rho_w)$ zakken. Het water, van gelijk volume, moet dezelfde diepte over hetzelfde oppervlak vullen, dus $r + s = \qty{574}{m}$,
$r\,\rho_m/(\rho_m - \rho_w) = 574$, en $r = 574\times2300/3300 =
\qty{400}{m}$: honderden meters, uit de spreidingssnelheid alleen.
\end{solution}

\begin{solution}{exo:b1:rifts-ridges-margins:11}
(a) $\sqrt{x - a} = \sqrt x\sqrt{1 - a/x} \approx \sqrt x\,(1 - a/2x)$ voor
$a \ll x$, dus
$\Delta d \approx C\sqrt{x/u}\cdot a/(2x) = Ca/(2\sqrt{ux})$.
(b) Bij $x = \qty{450}{km}$: exact \qty{1205}{m}, benaderd
$350\times350/(2\sqrt{10.6\times450}) = \qty{887}{m}$; bij \qty{1350}{km}:
\qty{550}{m} en \qty{512}{m}; bij \qty{3350}{km}: \qty{334}{m} en
\qty{325}{m}. De verhouding van de benaderde tot de exacte sprong hangt alleen af van $r = a/x$: $(r/2)/(1 - \sqrt{1 - r})$, die 0.9 is voor $r = 0.36$. De benadering is goed tot \qty{10}{\percent} vanaf $x = 350/0.36 \approx
\qty{970}{km}$, ongeveer drie verschuivingen van de rug.
\end{solution}

\begin{solution}{exo:b1:rifts-ridges-margins:12}
(a) $d_2 - d_1 = 3671.1 - 3606.8 = \qty{64.3}{m}$ bij \qty{10}{Myr};
$4072.8 - 4065.3 = \qty{7.5}{m}$ bij 20; $4415.4 - 4417.0 = \qty{-1.7}{m}$ bij
30; $4707.5 - 4713.6 = \qty{-6.1}{m}$ bij 40. (b) $d_2 - d_1$ is positief bij 20
en negatief bij 30, en beide wetten zijn continu: ze zijn daartussen gelijk.
Bisectie: bij 25, $d_2 - d_1 = +0.9$; bij 27.5, $-0.7$; bij 26.25, $+0.01$; bij
26.875, $-0.3$; bij 26.56, $-0.2$; bij 26.41, $-0.08$; bij 26.33, $-0.04$: het
nulpunt ligt tussen 26.25 en 26.33, bij \qty{26.3}{Myr}, waar beide
\qty{4294}{m} geven. (c) $d_1' = 175/\sqrt{26.3} = \qty{34.1}{m.Myr^{-1}}$ en
$d_2' = (3200/62.8)\,e^{-26.3/62.8} = \qty{33.5}{m.Myr^{-1}}$. De wetten hebben
daar bijna dezelfde waarde en dezelfde helling, dus de samengevoegde kromme
heeft geen zichtbare knik.
\end{solution}

\begin{solution}{pb:b1:rifts-ridges-margins:1}
\textbf{1.} \omterm{def:b1:cross-sections:throw}{Heave} $= $ \omterm{def:b1:cross-sections:throw}{worp}$/\tan$(helling): F1 \qty{346}{m}, F2 \qty{770}{m},
F3 \qty{2014}{m}, F4 \qty{163}{m}, F5 \qty{462}{m}, F6 \qty{1172}{m}. F3, met
de grootste \omterm{def:b1:cross-sections:throw}{worp} en \omterm{def:b1:cross-sections:throw}{heave}, is de randbreuk.

\textbf{2.} De uitrekking is de som, \qty{4927}{m}, ongeveer \qty{4.9}{km}. De
\omterm{def:b1:rifts-ridges-margins:rift}{slenk} was $60 - 4.9 = \qty{55.1}{km}$ breed vóór de breuken.

\textbf{3.} De normaal op een \omterm{def:b1:rifts-ridges-margins:rift}{slenk} die naar \ang{020} loopt, wijst naar
\ang{110}; de beweging, naar \ang{105}, maakt er \ang{5} mee. De loodrechte
component is $6.0\cos\ang{5} = \qty{5.98}{mm.yr^{-1}}$.

\textbf{4.} $\qty{4.927}{km}/\qty{5.98}{km.Myr^{-1}} = \qty{0.82}{Myr}$.

\textbf{5.} In \qty{3.2}{Myr} zijn de platen over de \omterm{def:b1:rifts-ridges-margins:rift}{slenk} met
$5.98\times3.2 = \qty{19.1}{km}$ uit elkaar bewogen. De breuken dragen
$4.93/19.1 = \qty{26}{\percent}$ ervan.

\textbf{6.} De overige driekwart, \qty{14.2}{km}, zit niet in de breuken van de
doorsnede. Het kan worden opgevangen door langs de as ingedrongen gangen, die de
\omterm{def:b1:rifts-ridges-margins:rift}{slenk} openen zonder dat er breuken in ontstaan, zoals in Afar, of door breuken
buiten de doorsnede. Met gangen van \qty{8}{m} breed zijn
$\qty{14200}{m}/\qty{8}{m} \approx 1780$ gangen nodig, één elke
$3.2\times10^6/1780 \approx 1800$ jaar gemiddeld.

\textbf{7.} Uit C2An: $(21.0 + 20.2)/2.581 = 41.2/2.581 =
\qty{15.96}{mm.yr^{-1}}$. Uit C3n: $(41.0 + 40.2)/5.235 = 81.2/5.235 =
\qty{15.51}{mm.yr^{-1}}$.

\textbf{8.} $(41.2 - 2.2)/2.581 = \qty{15.11}{mm.yr^{-1}}$ en
$(81.2 - 2.2)/5.235 = \qty{15.09}{mm.yr^{-1}}$: gecorrigeerd stemmen de twee
overeen tot op \qty{0.02}{mm.yr^{-1}}. De correctie is belangrijker voor de
jongere chron, waarvan de afstand korter is.

\textbf{9.} De verplaatsing naar buiten verschuift beide randen van een flank
met hetzelfde bedrag, dus die valt weg in hun verschil. Over beide flanken:
$(81.2 - 41.2)/(5.235 - 2.581) = 40.0/2.654 = \qty{15.07}{mm.yr^{-1}}$.

\textbf{10.} Noordoost: $(41.0 - 21.0)/2.654 = \qty{7.54}{mm.yr^{-1}}$; zuidwest:
$(40.2 - 20.2)/2.654 = \qty{7.54}{mm.yr^{-1}}$. De spreiding is symmetrisch; het
verschil van \qty{0.8}{km} tussen de flanken, voor beide chrons gelijk, is een
constante verschuiving, zoals wanneer de as \qty{0.4}{km} naast het midden ligt.

\textbf{11.} $\cos\Delta = \sin\ang{30.9}\sin\ang{18.0} +
\cos\ang{30.9}\cos\ang{18.0}\cos\ang{16.4} = 0.15870 + 0.78287 = 0.94157$, dus
$\Delta = \ang{19.69}$. Met $\omega R = 0.403\times(\pi/180)\times6371 =
\qty{44.81}{mm.yr^{-1}}$ is de snelheid $44.81\sin\ang{19.69} =
\qty{15.10}{mm.yr^{-1}}$, de gecorrigeerde snelheden van de vragen 8 en 9.

\textbf{12.} $\cos\Delta = \sin\ang{30.9}\sin\ang{26.0} +
\cos\ang{30.9}\cos\ang{26.0}\cos\ang{11.4} = 0.22512 + 0.75600 = 0.98112$, dus
$\Delta = \ang{11.15}$ en de snelheid is $44.81\sin\ang{11.15} =
\qty{8.7}{mm.yr^{-1}}$.

\textbf{13.} De voorspellingen, 15.1 en \qty{8.7}{mm.yr^{-1}}, komen overeen met
de gemeten 15 en \qty{9}{mm.yr^{-1}}. De snelheid is $\omega R\sin\Delta$, en de
pool ligt ten noordwesten van de Rode Zee, in het noordoosten van Libië: hoe
zuidelijker, hoe verder van de pool, hoe sneller de opening.

\textbf{14.} $15.09\times5.235 = \qty{79}{km}$ \omterm{def:b1:rifts-ridges-margins:oceanic-crust}{oceanische korst}.

\textbf{15.} $5000/120.6 = \qty{41}{km.Myr^{-1}}$, \qty{41}{mm.yr^{-1}} gemiddeld.

\textbf{16.} $5000/15.1 = \qty{331}{Myr}$ bij de centrale snelheid;
$5000/8.7 = \qty{575}{Myr}$ bij de noordelijke snelheid.

\textbf{17.} Bij opening rond een nabije pool wijdt de Rode Zee uit als een
schaar: haar breedte groeit evenredig met $\sin\Delta$, en zij blijft een wig
die naar de pool toe smaller wordt. Bij \qty{8.7}{mm.yr^{-1}} heeft het
noordelijke deel $79/8.7 = \qty{9.1}{Myr}$ nodig om dezelfde \qty{79}{km}
\omterm{def:b1:rifts-ridges-margins:oceanic-crust}{oceanische korst} te vormen die het centrale deel nu heeft.

\textbf{18.} Bij \qty{15}{mm.yr^{-1}} is de as van de Rode Zee een ultralangzame
rug; een as-dal tussen normale breuken wordt verwacht.

\textbf{19.} $2500 + 350\times1 = \qty{2850}{m}$;
$2500 + 350\sqrt{10} = \qty{3607}{m}$; $2500 + 350\sqrt{40} = \qty{4714}{m}$.

\textbf{20.} $t = (1500/350)^2 = \qty{18.4}{Myr}$.

\textbf{21.} In de eerste \qty{10}{Myr} zakt zij $350\sqrt{10} = \qty{1107}{m}$;
tussen 70 en \qty{80}{Myr}, $350(\sqrt{80} - \sqrt{70}) = 350\times(8.944 -
8.367) = \qty{202}{m}$, vijf keer minder in dezelfde tijd.

\textbf{22.} $\mathrm dd/\mathrm dt = C/(2\sqrt t) = 175/\sqrt{80} =
\qty{19.6}{m.Myr^{-1}}$, ongeveer \qty{0.02}{mm} per jaar.

\textbf{23.} $490/\sqrt{80} = \qty{54.8}{mW.m^{-2}}$.

\textbf{24.} Wet van de vierkantswortel: $2500 + 350\sqrt{80} = 2500 + 350\times8.944 =
\qty{5630}{m}$. Plaatmodel: $6400 - 3200\,e^{-80/62.8} = 6400 - 3200\times0.2798
= \qty{5505}{m}$. Ze verschillen \qty{126}{m}. Bij \qty{80}{Myr} gelden beide
wetten nog: de wet van de vierkantswortel tot ongeveer \qty{80}{Myr}, het
plaatmodel vanaf ongeveer \qty{20}{Myr}; daarna gaat de wet van de
vierkantswortel te diep.

\textbf{25.} De wet van de vierkantswortel geeft \qty{5630}{m}: op een tiende van
een kilometer nauwkeurig ligt de bodem van de nieuwe oceaan op een diepte van
\textbf{\qty[text-series-to-math=true]{5.6}{km}} wanneer zij 80 miljoen jaar oud
is (het plaatmodel geeft \qty{5.5}{km}).
\end{solution}
